\section*{الفصل \ref{ch:g11:forces-and-motion} --- القوى والحركة}

\begin{solution}{exo:g11:forces-and-motion:1}
(أ) السرعة ثابتة: \omterm{def:g11:forces-and-motion:netforce}{فالقوة المحصّلة} منعدمة. (ب) السرعة تتغير: فهي غير منعدمة،
معاكسة للحركة. (ج) الاتجاه يتغير (فالسرعة متجهة):
غير منعدمة، نحو داخل المنعطف. (د) السرعة ثابتة: منعدمة ---
إذ يتعادل \omterm{def:g10:universal-gravitation:weight}{الوزن} \omterm{def:g11:circuits-and-power:resistance}{والمقاومة}.
\end{solution}

\begin{solution}{exo:g11:forces-and-motion:2}
$d_1/d_2 = 1.0/3.0$ مع $d_1 + d_2 = \qty{0.80}{m}$:
أي $d_1 = \qty{0.20}{m}$ من الكرة \qty{3.0}{kg}. ولو تساوت الكتلتان: في
المنتصف، على بعد \qty{0.40}{m} من كلٍّ.
\end{solution}

\begin{solution}{exo:g11:forces-and-motion:3}
بالمقدار $F\Delta t$ نفسه وبثلاثة أضعاف الكتلة: يكون $\Delta v$ أصغر ثلاث مرات محمَّلة.
ولمضاهاة $\Delta v$ الخاص بالعربة الفارغة، ادفع أقوى ثلاث مرات (أو
أطول ثلاث مرات).
\end{solution}

\begin{solution}{exo:g11:forces-and-motion:4}
$\Delta \vect v$: أي \qty{2}{m/s}، موجَّهًا \emph{إلى الوراء} (من
\num{10} إلى \qty{8}{m/s}). وتتجه \omterm{def:g11:forces-and-motion:netforce}{القوة المحصّلة} في الجهة نفسها:
عكس $\vect v$ --- أي الحالة 2، التباطؤ.
\end{solution}

\begin{solution}{exo:g11:forces-and-motion:5}
(أ) سكوني، معاكس لدفعتك (فلا انزلاق بعد). (ب) حركي، معاكس
للانزلاق. (ج) سكوني، متجه \emph{صاعدًا} على المنحدر --- ضد
الانزلاق الذي تحاول \omterm{def:g11:fundamental-interactions:four}{الجاذبية} أن تبدأه.
\end{solution}

\begin{solution}{exo:g11:forces-and-motion:6}
$\Delta v \propto F\Delta t/m$: (أ) \qty{8.0}{m/s}؛ (ب) \qty{2.0}{m/s}؛
(ج) \qty{2.0}{m/s}؛ (د) $4.0 \times \frac{3 \times 2}{6} =
\qty{4.0}{m/s}$ --- فتتلاشى التغيرات الثلاثة.
\end{solution}

\begin{solution}{exo:g11:forces-and-motion:7}
$v' = \sqrt{8.0^2 + 6.0^2} = \qty{10.0}{m/s}$، بزاوية $\tan\theta =
6.0/8.0$، أي\ $\theta \approx 37^\circ$ شمال الشرق. فالسرعات تتجمع
\emph{متجهاتٍ}: إذ تتركّب الإسهامات المتعامدة بفيثاغورس،
لا بجمع الشدتين.
\end{solution}

\begin{solution}{exo:g11:forces-and-motion:8}
المثبَّت سلفًا: كتلة الكرة و$\Delta v$ الخاص بها (من السرعة الواردة إلى
الصفر)، ومن ثم الجداء $F\Delta t$. ومدّ $\Delta t$ من
\qty{0.02}{s} إلى \qty{0.08}{s} يقسم \omterm{def:g10:inertia:force}{القوة} المتوسطة على $4$.
\end{solution}

\begin{solution}{exo:g11:forces-and-motion:9}
\omterm{def:g11:forces-and-motion:friction}{الاحتكاك السكوني} بين المضمار والحذاء: فالقدم تدفع إلى الوراء
على الأرض، وتدفع \omterm{def:g10:inertia:contact}{قوة التماس} المتبادلة العدّاءَ
إلى الأمام. والمسامير ترفع الحدّ الأقصى \omterm{def:g10:inertia:inventory}{للاحتكاك} السكوني (فلا انزلاق عند بذل
الجهد كاملًا). وعلى جليد بلا \omterm{def:g10:inertia:inventory}{احتكاك} لا وجود لأي \omterm{def:g10:inertia:force}{قوة} أفقية خارجية، ومن ثم
لا تقدر سرعة \omterm{def:g11:forces-and-motion:com}{مركز الكتلة} على التغير --- فلا انطلاق ولا مشي.
\end{solution}

\begin{solution}{exo:g11:forces-and-motion:10}
نقطة تماس العجلة الدائرة لا تنزلق: \omterm{def:g10:inertia:inventory}{فالاحتكاك} سكوني،
وهو أقوى من \omterm{def:g11:forces-and-motion:friction}{الاحتكاك الحركي} في الانزلاق
(\cref{def:g11:forces-and-motion:friction}) --- ومن ثم توقف أقصر. أما \omterm{def:g11:forces-and-motion:friction}{الاحتكاك
الحركي} لإطار منزلق فيتجه عكس الانزلاق، مهما تكن
زاوية العجلات: فإدارة العجلة لا تغيّر شيئًا، ويضيع التوجيه.
\end{solution}

\begin{solution}{exo:g11:forces-and-motion:11}
\omterm{def:g10:inertia:inventory}{شدّ} السلك، متجهًا نحو الداخل إلى الرياضي. وبعد الإطلاق
تطير الكرة في خط مستقيم على المماسّ، بالسرعة
التي كانت لها (بالعطالة). أما ``سلك'' القمر فهو جذب الأرض
الثقالي؛ فأطفئ \omterm{def:g11:fundamental-interactions:four}{الجاذبية} يغادر على مماسّه، مستقيمًا
\omterm{def:g11:electric-gravitational-fields:capacitor}{منتظمًا}، إلى الأبد.
\end{solution}

\begin{solution}{exo:g11:forces-and-motion:12}
من الأعلى: $x = \dfrac{1.0 \times 0.60 + 2.0 \times 1.20}{3.0}
= \qty{1.00}{m}$. والمسافتان إلى مركزي الجزأين: \qty{0.40}{m}
(العصا) و\qty{0.20}{m} (الرأس)، والنسبة $2 : 1$ --- أي مقلوب
نسبة الكتلتين $1.0 : 2.0$، كما يقتضي \cref{def:g11:forces-and-motion:com}.
\end{solution}

\begin{solution}{exo:g11:forces-and-motion:13}
\emph{1.} بالمقدارين $F$ و$\Delta v$ نفسيهما وبضعف $m$: تحتاج الشاحنة ضعف
الزمن. \emph{2.} وبالسرعة المتوسطة نفسها في ضعف الزمن: ضعف
المسافة. \emph{3.} وبضعف $\Delta v$ عند $F$ و$m$ نفسيهما: ضعف
الزمن؛ لكن \omterm{def:g10:relative-motion:speed}{السرعة المتوسطة} تتضاعف هي أيضًا، فتُضرب المسافة
في $4$. فمضاعفة سرعتك تربّع مسافة
توقفك.
\end{solution}

\begin{solution}{exo:g11:forces-and-motion:14}
\emph{1.} $T = 27.3 \times \qty{86400}{s} \approx \qty{2.36e6}{s}$، ومنه
$v = \dfrac{2\pi \times \num{3.84e8}}{\num{2.36e6}} \approx
\qty{1.0e3}{m/s} \approx \qty{1}{km/s}$. \emph{2.} نحو الأرض
(إلى الداخل)؛ وتوفّره جاذبية الأرض. \emph{3.} فالقمر \emph{يسقط}
فعلًا: إذ يتجه تغير سرعته نحو الأرض في كل لحظة. أما
سرعته المماسّية فلا تفعل إلا أن تضمن أنه يظلّ يخطئ الهدف --- فهو يسقط
\emph{حول} الأرض بدل أن يسقط فيها.
\end{solution}

\begin{solution}{exo:g11:forces-and-motion:15}
بالمقدارين $m$ و$\Delta v$ نفسيهما، يكون $F\Delta t$ ثابتًا: وكون $\Delta t$ أطول خمسين
مرة على الرغوة يعني \omterm{def:g10:inertia:force}{قوة} متوسطة أصغر خمسين مرة.
وثني الركبتين يمدّ $\Delta t$ الخاص بالهبوط من هزّة إلى
ثنية طويلة؛ وحصيرة السقوط، ميكانيكيًّا، مصنوعة من أجزاء إضافية من
الألف من \omterm{def:g10:orders-of-magnitude:unit}{الثانية}.
\end{solution}

\begin{solution}{pb:g11:forces-and-motion:1}
\textbf{1.} $50/3.6 \approx \qty{13.9}{m/s}$ (ولنسمّها \qty{14}{m/s}).
\textbf{2.} $\Delta v \approx \qty{14}{m/s}$: فالراكب يتحرك بسرعة
\qty{14}{m/s} قبله وبسرعة صفر بعده --- والطرفان يثبّتهما
التصادم، فلا جهاز يقدر أن يمسّ $\Delta v$.
\textbf{3.} $\Delta v \propto F\Delta t/m$
(\cref{thm:g11:forces-and-motion:secondlaw}) مع تثبيت $\Delta v$ و$m$
يفرض أن يكون $F\Delta t$ ثابتًا. فمقبض المصمّم الوحيد هو
$\Delta t$.
\textbf{4.} $0.150/0.050 = 3$: فمنطقة الانسحاق تقسم \omterm{def:g10:inertia:force}{القوة}
المتوسطة على $3$.
\textbf{5.} إلى الوراء، عكس الحركة --- فالسرعة تتقلص من
\qty{14}{m/s} إلى الصفر.
\textbf{6.} المقدمة القابلة للطيّ آلة تجعل التوقف يدوم
أطول: فالمعدن الذي ينسحق يشتري أجزاءً من الألف من \omterm{def:g10:orders-of-magnitude:unit}{الثانية}.
\textbf{7.} بالمقدارين $\Delta v$ و$m$ نفسيهما، ومع $\Delta t$ أقصر ثلاث مرات:
تتلقى الدمية القديمة ثلاثة أضعاف \omterm{def:g10:inertia:force}{القوة}.
\textbf{8.} \omterm{def:g10:relative-motion:speed}{السرعة المتوسطة} $\approx \qty{6.9}{m/s}$؛ والمسافة $\approx
6.9 \times 0.150 \approx \qty{1.0}{m}$ --- أي تقريبًا \omterm{def:g10:orders-of-magnitude:unit}{المتر} من
السيارة القابلة للتشوّه الذي توفّره منطقة انسحاق \omterm{def:g11:lenses-and-eye:images}{حقيقية}.
\textbf{9.} مقدمة طولها \qty{10}{m} لا تُقاد ولا تُركن، وكتلتها
هي نفسها كانت ستفاقم كل تصادم. أما حجرة الركاب فيجب ألّا
تنسحق: فهي تحفظ حيّز النجاة.
\textbf{10.} ``\dots فعليك أن تختار $\Delta t$.''
\textbf{11.} $0.150/0.005 = 30$: فالتوقف على الزجاج الأمامي أعنف
ثلاثين مرة من التوقف بالحزام.
\textbf{12.} الحزام الذي يتمطط قليلًا يطيل $\Delta t$ الخاص بالراكب
قليلًا أكثر (ويوزّع \omterm{def:g10:inertia:force}{القوة})؛ أما طقم فولاذي
فكان سيفرض أقصر توقف لهيكل السيارة على ألين الركاب.
\textbf{13.} $0.080/0.010 = 8$: فالوسادة الهوائية تقسم \omterm{def:g10:inertia:force}{قوة} الرأس على
$8$.
\textbf{14.} بالمقدارين $\Delta v$ و$\Delta t$ نفسيهما: تتغير \omterm{def:g10:inertia:force}{القوة} مع
الكتلة، فتكون أكبر $70/20 = 3.5$ مرة للبالغ.
\textbf{15.} التغير المطلوب: $14/0.150 \approx \qty{93}{m/s}$ في كل
\omterm{def:g10:orders-of-magnitude:unit}{ثانية} --- أي نحو $9.5$ ضعف ما تنتجه \omterm{def:g11:fundamental-interactions:four}{الجاذبية}
(\qty{9.8}{m/s} في \omterm{def:g10:orders-of-magnitude:unit}{الثانية}). فيجب أن تشدّ الذراعان نحو $9.5$ ضعف
\omterm{def:g10:universal-gravitation:weight}{وزن} الرضيع: أي كأنك تحمل حِملًا قدره \qty{95}{kg}. ولا ذراعين تقدران؛ فيطير الرضيع
إلى الأمام.
\textbf{16.} حزام البالغ يوجّه \omterm{def:g10:inertia:force}{القوة} إلى الوركين والصدر؛ وعلى
الطفل كان سيحمّل البطن والرقبة. أما مقعد الطفل فيطبّق
\omterm{def:g10:inertia:force}{القوة} (الأصغر) على أجزاء قوية من الجسم، وتمدّ حشوته
$\Delta t$ أكثر.
\textbf{17.} إن تشوّهت الحجرة فقد الركاب حيّز
نجاتهم وصدمهم الهيكل نفسه. ولا ينبغي أن يصطدم شيء
بالحجرة الصلبة مباشرة: فتوقفات الركاب يتولاها الحزام و
الوسادة الهوائية (بزمن $\Delta t$ طويل)، بينما ترسي الحجرة الطرفين المنسحقين.
\textbf{18.} بالتناظر تتوقف السيارتان عند مستوي التماس: فيخضع كل
سائق للتغير $\Delta v \approx \qty{14}{m/s}$ على مدى توقف بمنطقة
انسحاق --- أي اختبار الجدار أساسًا، لا ``تصادمًا مضاعفًا''.
\textbf{19.} القوتان متساويتان متعاكستان، لكن
$\Delta v \propto F\Delta t/m$: فالسيارة الصغيرة، ذات $m$ الصغيرة،
تعاني تغير السرعة الكبير.
\textbf{20.} منطقة الانسحاق، وحزام المقعد المتمطط، والوسادة الهوائية، وحشوة
مقعد الطفل --- وكلٌّ منها يمدّ المقدار نفسه: $\Delta t$.
\textbf{21.} يثبّت التصادم الكتلة وتغير السرعة، ومن ثم
جداء \omterm{def:g10:inertia:force}{القوة} $\times$ الزمن: فالتصادم اللطيف تصادم طويل ---
فمدّ الزمن تهبط \omterm{def:g10:inertia:force}{القوة} بالنسبة نفسها.
\end{solution}
